\documentclass[12pt]{article}
\usepackage[utf8]{inputenc}
\usepackage[T1]{fontenc}
\usepackage{hyperref}
\title{Northern Muji Language Resources for K-12 Curriculum Development}
\author{Compiled by Claude Sonnet 5.5 (claude-sonnet-5-5) for Andr\'as Kornai}
\date{}
\begin{document}
\maketitle

\section{Introduction}
Northern Muji (ISO 639-3 ymx; Glottocode nort2716) is a Sino-Tibetan (Loloish) language of Yunnan, China. Alternate names include Bokha, Hlaka Mujima and Phula. The name ``Muji'' is shared by several related languages (Southern Muji, Qila Muji, Muzi and others), so resources must be checked carefully. Speaker estimates differ: about 9,000 (Verbix) and 15,000 for the people group (Joshua Project). The search found no school materials, primers or native-language curriculum for this language.

\section{Primary Educational Resources}
No primers, textbooks or school programmes for Northern Muji were found in search results. The language is not known to be taught in schools.

\section{Linguistic and Reference Materials}
\begin{itemize}
\item Wikipedia, overview of the Muji language cluster: \url{https://en.wikipedia.org/wiki/Muji_language}
\item Verbix Endangered Languages, Northern Muji: \url{https://docs.verbix.com/EndangeredLanguages/NorthernMuji}
\item Endangered Languages Project, bibliography page for Northern Muji: \url{https://endangeredlanguages.com/lang/3175/bibliography}
\end{itemize}

\section{Religious and Community Resources}
\begin{itemize}
\item Joshua Project, Northern Muji people group profile (demographics and location in southern Yunnan): \url{https://joshuaproject.net/people_groups/20606}
\item Joshua Project, language page for ymx: \url{https://joshuaproject.net/languages/ymx}
\end{itemize}
No scripture translations or community organisations were identified.

\section{Digital and Audio Resources}
No audio, video, corpus or software resources were found in the search results.

\section{Curriculum Development Context}
Little published material exists. Any curriculum would need to start with documentation work, such as fieldwork, orthography development and recordings, together with the community. The Endangered Languages Project bibliography page above is the best lead for locating academic descriptions. Care is needed to separate Northern Muji from the other Muji varieties.

\end{document}